\documentclass[UTF8]{ctexart}
\usepackage{ctex}
\usepackage{amsmath,amssymb,amsthm,mathtools,bm}
\usepackage{geometry}
\geometry{a4paper,margin=1in}

\newtheorem{assumption}{假设}
\newtheorem{definition}{定义}
\newtheorem{lemma}{引理}
\newtheorem{proposition}{命题}
\newtheorem{theorem}{定理}
\newtheorem{remark}{注}

\DeclareMathOperator{\spanop}{span}
\DeclareMathOperator{\dist}{dist}

\begin{document}
%============================================================
% Fourier--GLOF post-processing: weighted L2 error estimate
%============================================================

% Suggested packages:
% \usepackage{amsmath,amssymb,amsthm}
%
% Suggested theorem environments:
% \newtheorem{theorem}{Theorem}
% \newtheorem{lemma}{Lemma}
% \newtheorem{proposition}{Proposition}
% \newtheorem{assumption}{Assumption}
% \newtheorem{remark}{Remark}

\section{A weighted $L^2$ estimate for Fourier--GLOF post-processing}

Throughout this section, let
\[
    I=(0,1),\qquad L(x):=-\log x .
\]
For $\alpha>-1$ and $\lambda>-1$, define the weighted inner product
\[
    (u,v)_{\alpha,\lambda}
    :=
    \int_0^1 u(x)\overline{v(x)}
    \chi_{\alpha,\lambda}(x)\,dx,
    \qquad
    \chi_{\alpha,\lambda}(x):=L(x)^\alpha x^\lambda,
\]
and the corresponding norm
\[
    \|u\|_{\alpha,\lambda}
    :=
    (u,u)_{\alpha,\lambda}^{1/2}.
\]

Let $\alpha,\beta>-1$ and $\lambda>-1$. The generalized log orthogonal functions
are defined by
\[
    \mathcal S_n^{(\alpha,\beta,\lambda)}(x)
    :=
    x^{\frac{\beta-\lambda}{2}}
    L_n^{(\alpha)}\bigl((\beta+1)L(x)\bigr),
    \qquad n=0,1,2,\dots,
\]
where $L_n^{(\alpha)}$ is the Laguerre polynomial of degree $n$.
These functions satisfy
\[
    \int_0^1
    \mathcal S_n^{(\alpha,\beta,\lambda)}(x)
    \mathcal S_m^{(\alpha,\beta,\lambda)}(x)
    \chi_{\alpha,\lambda}(x)\,dx
    =
    \gamma_n^{(\alpha,\beta)}\delta_{mn},
\]
where
\[
    \gamma_n^{(\alpha,\beta)}
    =
    \frac{\Gamma(n+\alpha+1)}
    {(\beta+1)^{\alpha+1}\Gamma(n+1)}.
\]

Let
\[
    V_M^{\alpha,\beta,\lambda}
    :=
    \operatorname{span}
    \left\{
    \mathcal S_0^{(\alpha,\beta,\lambda)},
    \mathcal S_1^{(\alpha,\beta,\lambda)},
    \dots,
    \mathcal S_M^{(\alpha,\beta,\lambda)}
    \right\}.
\]
The weighted $L^2$-orthogonal projection
\[
    \pi_M^{\alpha,\beta,\lambda}
    :
    L^2_{\chi_{\alpha,\lambda}}(0,1)
    \longrightarrow
    V_M^{\alpha,\beta,\lambda}
\]
is given by
\[
    \pi_M^{\alpha,\beta,\lambda}u
    =
    \sum_{n=0}^M
    \widehat u_n^{\alpha,\beta,\lambda}
    \mathcal S_n^{(\alpha,\beta,\lambda)},
\]
where
\[
    \widehat u_n^{\alpha,\beta,\lambda}
    =
    \frac{1}{\gamma_n^{(\alpha,\beta)}}
    \int_0^1
    u(x)
    \mathcal S_n^{(\alpha,\beta,\lambda)}(x)
    \chi_{\alpha,\lambda}(x)\,dx .
\]

We consider
\[
    f(x)=a(x)+b(x)x^rL(x)^k,
    \qquad r>0,\qquad k\in\mathbb N.
\]
Let the Fourier partial sum of $f$ be
\[
    f_N^{\rm F}(x)
    =
    \sum_{|\ell|\le N}c_\ell e^{2\pi i\ell x},
    \qquad
    c_\ell
    =
    \int_0^1 f(x)e^{-2\pi i\ell x}\,dx .
\]

%------------------------------------------------------------
\subsection{Assumptions}

\begin{assumption}[Basic regularity of the data]
\label{ass:data}
Assume that $a,b$ are analytic in a neighbourhood of $[0,1]$.
In particular,
\[
    f\in W^{1,1}(0,1).
\]
Define
\[
    B_f
    :=
    \frac{
    |f(1)-f(0)|+\|f'\|_{L^1(0,1)}
    }{2\pi}.
\]
\end{assumption}

Under Assumption~\ref{ass:data}, integration by parts gives, for every
$\ell\ne 0$,
\[
    |c_\ell|
    \le
    \frac{B_f}{|\ell|}.
\]

\begin{assumption}[Weighted $L^2$ GLOF regularization]
\label{ass:regularization}
Assume that there exists a known sequence $\mathcal R_M$ such that
\[
    \|f-\pi_M^{\alpha,\beta,\lambda}f\|_{\alpha,\lambda}
    \le
    \mathcal R_M.
\]
In applications one typically aims to prove
\[
    \mathcal R_M
    \le
    C_{\rm reg} M^{s_{\rm reg}}q_{\rm reg}^M,
    \qquad
    0<q_{\rm reg}<1.
\]
\end{assumption}

The role of Assumption~\ref{ass:regularization} is to isolate the
regularization part of the proof. It is the part supplied by the
GLOF approximation theory. For example, it is verified for finite
linear combinations of the model singular functions
$x^p(-\log x)^j$ by the single-term GLOF estimate stated below.

%------------------------------------------------------------
\subsection{A sufficient condition for the regularization estimate}

\begin{lemma}[Single-term GLOF projection estimate]
\label{lem:single-term}
Let
\[
    g_{p,j}(x):=x^pL(x)^j,
    \qquad p\ge 0,\qquad j\in\mathbb N_0.
\]
Assume
\[
    \beta>\lambda,
    \qquad
    \lambda>-1-2p.
\]
Define
\[
    R_p^{\alpha,\beta,\lambda}
    :=
    \left|
    \frac{2p+\lambda-\beta}
    {2p+2+\lambda+\beta}
    \right|.
\]
Then
\[
    R_p^{\alpha,\beta,\lambda}<1,
\]
and for all sufficiently large $M$,
\[
    \left\|
    g_{p,j}
    -
    \pi_M^{\alpha,\beta,\lambda}g_{p,j}
    \right\|_{\alpha,\lambda}
    \le
    C_{p,j,\alpha,\beta,\lambda}
    M^{\frac{\alpha+1}{2}+j}
    \left(R_p^{\alpha,\beta,\lambda}\right)^M .
\]
\end{lemma}

\begin{proposition}[Regularization for finite singular expansions]
\label{prop:finite-reg}
Suppose that
\[
    a(x)=\sum_{j=0}^{J_a}a_jx^j,
    \qquad
    b(x)=\sum_{j=0}^{J_b}b_jx^j .
\]
Assume
\[
    \beta>\lambda,
    \qquad
    \lambda>-1.
\]
Define the finite set of exponents
\[
    \mathcal P
    :=
    \{0,1,\dots,J_a\}
    \cup
    \{r,r+1,\dots,r+J_b\},
\]
and
\[
    \rho_{\rm reg}
    :=
    \max_{p\in\mathcal P}
    R_p^{\alpha,\beta,\lambda},
    \qquad
    R_p^{\alpha,\beta,\lambda}
    :=
    \left|
    \frac{2p+\lambda-\beta}
    {2p+2+\lambda+\beta}
    \right|.
\]
Then
\[
    0<\rho_{\rm reg}<1,
\]
and for all sufficiently large $M$,
\[
    \left\|
    f-\pi_M^{\alpha,\beta,\lambda}f
    \right\|_{\alpha,\lambda}
    \le
    C_{\rm reg}
    M^{\frac{\alpha+1}{2}+k}
    \rho_{\rm reg}^M .
\]
\end{proposition}

\begin{proof}
Since
\[
    f(x)
    =
    \sum_{j=0}^{J_a}a_jx^j
    +
    \sum_{j=0}^{J_b}b_jx^{r+j}L(x)^k,
\]
we have, by the triangle inequality,
\[
\begin{aligned}
    \left\|
    f-\pi_M^{\alpha,\beta,\lambda}f
    \right\|_{\alpha,\lambda}
    &\le
    \sum_{j=0}^{J_a}|a_j|
    \left\|
    x^j-\pi_M^{\alpha,\beta,\lambda}x^j
    \right\|_{\alpha,\lambda}
    \\
    &\quad+
    \sum_{j=0}^{J_b}|b_j|
    \left\|
    x^{r+j}L(x)^k
    -
    \pi_M^{\alpha,\beta,\lambda}
    \bigl(x^{r+j}L(x)^k\bigr)
    \right\|_{\alpha,\lambda}.
\end{aligned}
\]
Applying Lemma~\ref{lem:single-term} to each term gives
\[
    \left\|
    f-\pi_M^{\alpha,\beta,\lambda}f
    \right\|_{\alpha,\lambda}
    \le
    C_{\rm reg}
    M^{\frac{\alpha+1}{2}+k}
    \rho_{\rm reg}^M .
\]
Since the set $\mathcal P$ is finite and $\beta>\lambda$, one has
\[
    R_p^{\alpha,\beta,\lambda}<1
    \qquad\text{for each }p\in\mathcal P,
\]
hence $\rho_{\rm reg}<1$.
\end{proof}

%------------------------------------------------------------
\subsection{The Fourier truncation term}

For the truncation estimate, define
\[
    \eta:=\frac{\beta+\lambda}{2}.
\]
For each $n\ge 0$, define the GLOF testing kernel
\[
    W_n^{\alpha,\beta,\lambda}(x)
    :=
    \mathcal S_n^{(\alpha,\beta,\lambda)}(x)
    \chi_{\alpha,\lambda}(x).
\]
Equivalently,
\[
    W_n^{\alpha,\beta,\lambda}(x)
    =
    x^\eta
    L(x)^\alpha
    L_n^{(\alpha)}\bigl((\beta+1)L(x)\bigr).
\]

Let $d\in\mathbb N$ satisfy
\[
    d<1+\min\{\alpha,\eta\}.
\]
Define
\[
    A_{n,d}^{\alpha,\beta,\lambda}
    :=
    \left\|
    \frac{d^d}{dx^d}
    W_n^{\alpha,\beta,\lambda}
    \right\|_{L^1(0,1)},
\]
and
\[
    \mathcal K_{M,d}^{\alpha,\beta,\lambda}
    :=
    \left(
    \sum_{n=0}^M
    \frac{
    \left(A_{n,d}^{\alpha,\beta,\lambda}\right)^2
    }{
    \gamma_n^{(\alpha,\beta)}
    }
    \right)^{1/2}.
\]

%------------------------------------------------------------
\begin{lemma}[Oscillatory kernel estimate]
\label{lem:kernel-estimate}
Let
\[
    K_{\ell n}
    :=
    \int_0^1
    e^{2\pi i\ell x}
    W_n^{\alpha,\beta,\lambda}(x)\,dx,
    \qquad
    \ell\ne 0.
\]
If
\[
    d<1+\min\{\alpha,\eta\},
    \qquad
    \eta=\frac{\beta+\lambda}{2},
\]
then
\[
    |K_{\ell n}|
    \le
    \frac{
    A_{n,d}^{\alpha,\beta,\lambda}
    }{
    (2\pi|\ell|)^d
    }.
\]
\end{lemma}

\begin{proof}
Near $x=0$,
\[
    W_n^{\alpha,\beta,\lambda}(x)
    =
    O\bigl(x^\eta L(x)^{\alpha+n}\bigr),
\]
and near $x=1$,
\[
    W_n^{\alpha,\beta,\lambda}(x)
    =
    O\bigl((1-x)^\alpha\bigr).
\]
More generally, for $0\le j\le d-1$,
\[
    \frac{d^j}{dx^j}
    W_n^{\alpha,\beta,\lambda}(x)
    \to 0
    \qquad
    \text{as }x\to 0^+\text{ and }x\to 1^-,
\]
because
\[
    j\le d-1<\eta,
    \qquad
    j\le d-1<\alpha.
\]
Therefore integration by parts $d$ times produces no boundary terms:
\[
    K_{\ell n}
    =
    \frac{(-1)^d}{(2\pi i\ell)^d}
    \int_0^1
    e^{2\pi i\ell x}
    \frac{d^d}{dx^d}
    W_n^{\alpha,\beta,\lambda}(x)\,dx .
\]
Taking absolute values yields
\[
    |K_{\ell n}|
    \le
    \frac{1}{(2\pi|\ell|)^d}
    \left\|
    \frac{d^d}{dx^d}
    W_n^{\alpha,\beta,\lambda}
    \right\|_{L^1(0,1)}
    =
    \frac{
    A_{n,d}^{\alpha,\beta,\lambda}
    }{
    (2\pi|\ell|)^d
    }.
\]
\end{proof}

%------------------------------------------------------------
\begin{lemma}[A computable upper bound for $A_{n,d}^{\alpha,\beta,\lambda}$]
\label{lem:A-bound}
Let
\[
    b_{n,j}^{(\alpha)}
    :=
    \frac{(-1)^j}{j!}
    \binom{n+\alpha}{n-j},
    \qquad
    0\le j\le n.
\]
Define coefficients $P_{d,h}(\eta,s)$ recursively by
\[
    P_{0,0}(\eta,s)=1,
\]
\[
    P_{d,h}(\eta,s)=0
    \qquad
    \text{if }h<0\text{ or }h>d,
\]
and
\[
    P_{d+1,h}(\eta,s)
    =
    (\eta-d)P_{d,h}(\eta,s)
    -
    (s-h+1)P_{d,h-1}(\eta,s).
\]
Then
\[
\begin{aligned}
    A_{n,d}^{\alpha,\beta,\lambda}
    \le
    \mathfrak A_{n,d}^{\alpha,\beta,\lambda},
\end{aligned}
\]
where
\[
\boxed{
\begin{aligned}
    \mathfrak A_{n,d}^{\alpha,\beta,\lambda}
    &:=
    \sum_{j=0}^n
    \left|b_{n,j}^{(\alpha)}\right|
    (\beta+1)^j
    \sum_{h=0}^d
    \left|P_{d,h}(\eta,\alpha+j)\right|
    \frac{
    \Gamma(\alpha+j-h+1)
    }{
    (\eta-d+1)^{\alpha+j-h+1}
    } .
\end{aligned}
}
\]
Consequently,
\[
    \mathcal K_{M,d}^{\alpha,\beta,\lambda}
    \le
    \left(
    \sum_{n=0}^M
    \frac{
    \left(
    \mathfrak A_{n,d}^{\alpha,\beta,\lambda}
    \right)^2
    }{
    \gamma_n^{(\alpha,\beta)}
    }
    \right)^{1/2}.
\]
\end{lemma}

\begin{proof}
Using the Laguerre expansion
\[
    L_n^{(\alpha)}(z)
    =
    \sum_{j=0}^n
    \frac{(-1)^j}{j!}
    \binom{n+\alpha}{n-j}
    z^j,
\]
we obtain
\[
    W_n^{\alpha,\beta,\lambda}(x)
    =
    \sum_{j=0}^n
    b_{n,j}^{(\alpha)}
    (\beta+1)^j
    x^\eta L(x)^{\alpha+j}.
\]
A direct induction gives
\[
    \frac{d^d}{dx^d}
    \bigl[x^\eta L(x)^s\bigr]
    =
    x^{\eta-d}
    \sum_{h=0}^d
    P_{d,h}(\eta,s)L(x)^{s-h}.
\]
Therefore
\[
\begin{aligned}
    A_{n,d}^{\alpha,\beta,\lambda}
    &\le
    \sum_{j=0}^n
    \left|b_{n,j}^{(\alpha)}\right|
    (\beta+1)^j
    \sum_{h=0}^d
    \left|P_{d,h}(\eta,\alpha+j)\right|
    \int_0^1
    x^{\eta-d}
    L(x)^{\alpha+j-h}\,dx .
\end{aligned}
\]
Since
\[
    d<\eta+1,
    \qquad
    d<\alpha+1,
\]
the integral is finite and equals
\[
    \int_0^1
    x^{\eta-d}
    L(x)^{\alpha+j-h}\,dx
    =
    \frac{
    \Gamma(\alpha+j-h+1)
    }{
    (\eta-d+1)^{\alpha+j-h+1}
    }.
\]
Substitution gives the claimed bound.
\end{proof}

%------------------------------------------------------------
\subsection{Main estimate}

\begin{theorem}[Weighted $L^2$ Fourier--GLOF post-processing error]
\label{thm:main-L2}
Let Assumption~\ref{ass:data} hold. Let
\[
    \alpha,\beta>-1,
    \qquad
    \lambda>-1,
    \qquad
    \eta=\frac{\beta+\lambda}{2}.
\]
Let $d\in\mathbb N$ satisfy
\[
    d<1+\min\{\alpha,\eta\}.
\]
Then
\[
\boxed{
\begin{aligned}
    &
    \left\|
    f-\pi_M^{\alpha,\beta,\lambda}
    \bigl(f_N^{\rm F}\bigr)
    \right\|_{\alpha,\lambda}^2
    \\
    &\qquad
    =
    \left\|
    f-\pi_M^{\alpha,\beta,\lambda}f
    \right\|_{\alpha,\lambda}^2
    +
    \left\|
    \pi_M^{\alpha,\beta,\lambda}
    \bigl(f-f_N^{\rm F}\bigr)
    \right\|_{\alpha,\lambda}^2 .
\end{aligned}
}
\]
Moreover,
\[
\boxed{
    \left\|
    \pi_M^{\alpha,\beta,\lambda}
    \bigl(f-f_N^{\rm F}\bigr)
    \right\|_{\alpha,\lambda}
    \le
    \frac{
    2B_f
    }{
    d(2\pi)^d
    }
    N^{-d}
    \mathcal K_{M,d}^{\alpha,\beta,\lambda}.
}
\]
Consequently, if Assumption~\ref{ass:regularization} holds, then
\[
\boxed{
\begin{aligned}
    &
    \left\|
    f-\pi_M^{\alpha,\beta,\lambda}
    \bigl(f_N^{\rm F}\bigr)
    \right\|_{\alpha,\lambda}
    \\
    &\qquad
    \le
    \mathcal R_M
    +
    \frac{
    2B_f
    }{
    d(2\pi)^d
    }
    N^{-d}
    \mathcal K_{M,d}^{\alpha,\beta,\lambda}.
\end{aligned}
}
\]
If, in addition, the finite-expansion assumption of
Proposition~\ref{prop:finite-reg} holds, then
\[
\boxed{
\begin{aligned}
    &
    \left\|
    f-\pi_M^{\alpha,\beta,\lambda}
    \bigl(f_N^{\rm F}\bigr)
    \right\|_{\alpha,\lambda}
    \\
    &\qquad
    \le
    C_{\rm reg}
    M^{\frac{\alpha+1}{2}+k}
    \rho_{\rm reg}^M
    +
    \frac{
    2B_f
    }{
    d(2\pi)^d
    }
    N^{-d}
    \mathcal K_{M,d}^{\alpha,\beta,\lambda}.
\end{aligned}
}
\]
\end{theorem}

\begin{proof}
Let
\[
    V_M
    :=
    V_M^{\alpha,\beta,\lambda}.
\]
Since
\[
    \pi_M^{\alpha,\beta,\lambda}f\in V_M,
\]
and $\pi_M^{\alpha,\beta,\lambda}$ is the weighted $L^2$-orthogonal projection,
we have
\[
    f-\pi_M^{\alpha,\beta,\lambda}f
    \perp_{\alpha,\lambda}
    V_M.
\]
On the other hand,
\[
    \pi_M^{\alpha,\beta,\lambda}
    \bigl(f-f_N^{\rm F}\bigr)
    \in V_M.
\]
Therefore
\[
    f-\pi_M^{\alpha,\beta,\lambda}
    \bigl(f_N^{\rm F}\bigr)
    =
    \bigl(f-\pi_M^{\alpha,\beta,\lambda}f\bigr)
    +
    \pi_M^{\alpha,\beta,\lambda}
    \bigl(f-f_N^{\rm F}\bigr)
\]
is an orthogonal decomposition in $L^2_{\chi_{\alpha,\lambda}}(0,1)$.
Hence the Pythagorean identity follows:
\[
\begin{aligned}
    &
    \left\|
    f-\pi_M^{\alpha,\beta,\lambda}
    \bigl(f_N^{\rm F}\bigr)
    \right\|_{\alpha,\lambda}^2
    \\
    &\qquad
    =
    \left\|
    f-\pi_M^{\alpha,\beta,\lambda}f
    \right\|_{\alpha,\lambda}^2
    +
    \left\|
    \pi_M^{\alpha,\beta,\lambda}
    \bigl(f-f_N^{\rm F}\bigr)
    \right\|_{\alpha,\lambda}^2 .
\end{aligned}
\]

It remains to estimate the second term. Set
\[
    e_N(x):=f(x)-f_N^{\rm F}(x).
\]
Then
\[
    \pi_M^{\alpha,\beta,\lambda}e_N
    =
    \sum_{n=0}^M
    \varepsilon_{n,N}
    \mathcal S_n^{(\alpha,\beta,\lambda)}(x),
\]
where
\[
    \varepsilon_{n,N}
    =
    \frac{1}{\gamma_n^{(\alpha,\beta)}}
    \int_0^1
    e_N(x)
    \mathcal S_n^{(\alpha,\beta,\lambda)}(x)
    \chi_{\alpha,\lambda}(x)\,dx .
\]
Using
\[
    e_N(x)
    =
    \sum_{|\ell|>N}c_\ell e^{2\pi i\ell x}
\]
in the $L^2$ sense, we obtain
\[
    \varepsilon_{n,N}
    =
    \frac{1}{\gamma_n^{(\alpha,\beta)}}
    \sum_{|\ell|>N}
    c_\ell K_{\ell n},
\]
with
\[
    K_{\ell n}
    =
    \int_0^1
    e^{2\pi i\ell x}
    W_n^{\alpha,\beta,\lambda}(x)\,dx .
\]
By Lemma~\ref{lem:kernel-estimate},
\[
    |K_{\ell n}|
    \le
    \frac{
    A_{n,d}^{\alpha,\beta,\lambda}
    }{
    (2\pi|\ell|)^d
    }.
\]
By Assumption~\ref{ass:data},
\[
    |c_\ell|\le \frac{B_f}{|\ell|},
    \qquad \ell\ne 0.
\]
Thus
\[
\begin{aligned}
    |\varepsilon_{n,N}|
    &\le
    \frac{1}{\gamma_n^{(\alpha,\beta)}}
    \sum_{|\ell|>N}
    |c_\ell|
    |K_{\ell n}|
    \\
    &\le
    \frac{
    B_f A_{n,d}^{\alpha,\beta,\lambda}
    }{
    (2\pi)^d
    \gamma_n^{(\alpha,\beta)}
    }
    \sum_{|\ell|>N}
    |\ell|^{-d-1}.
\end{aligned}
\]
Since
\[
    \sum_{|\ell|>N}
    |\ell|^{-d-1}
    =
    2\sum_{\ell=N+1}^\infty \ell^{-d-1}
    \le
    \frac{2}{d}N^{-d},
\]
we have
\[
    |\varepsilon_{n,N}|
    \le
    \frac{
    2B_f
    }{
    d(2\pi)^d
    }
    N^{-d}
    \frac{
    A_{n,d}^{\alpha,\beta,\lambda}
    }{
    \gamma_n^{(\alpha,\beta)}
    }.
\]
Using the GLOF orthogonality,
\[
\begin{aligned}
    \left\|
    \pi_M^{\alpha,\beta,\lambda}e_N
    \right\|_{\alpha,\lambda}^2
    &=
    \sum_{n=0}^M
    \gamma_n^{(\alpha,\beta)}
    |\varepsilon_{n,N}|^2
    \\
    &\le
    \left(
    \frac{
    2B_f
    }{
    d(2\pi)^d
    }
    N^{-d}
    \right)^2
    \sum_{n=0}^M
    \frac{
    \left(A_{n,d}^{\alpha,\beta,\lambda}\right)^2
    }{
    \gamma_n^{(\alpha,\beta)}
    }.
\end{aligned}
\]
Therefore
\[
    \left\|
    \pi_M^{\alpha,\beta,\lambda}
    \bigl(f-f_N^{\rm F}\bigr)
    \right\|_{\alpha,\lambda}
    \le
    \frac{
    2B_f
    }{
    d(2\pi)^d
    }
    N^{-d}
    \mathcal K_{M,d}^{\alpha,\beta,\lambda}.
\]
Combining this truncation estimate with
Assumption~\ref{ass:regularization} gives the general bound.
If Proposition~\ref{prop:finite-reg} is applicable, substituting its estimate for
$\mathcal R_M$ gives the final stated inequality.
\end{proof}

%------------------------------------------------------------
\subsection{Parameter-dependent form}

The previous theorem keeps all parameters explicit. Therefore, for any parameter choice
\[
    M=M_N,
    \qquad
    \alpha=\alpha_N,
    \qquad
    \beta=\beta_N,
    \qquad
    \lambda=\lambda_N,
\]
and for any integer
\[
    d_N<1+\min\left\{
    \alpha_N,\frac{\beta_N+\lambda_N}{2}
    \right\},
\]
one obtains
\[
\boxed{
\begin{aligned}
    &
    \left\|
    f-\pi_{M_N}^{\alpha_N,\beta_N,\lambda_N}
    \bigl(f_N^{\rm F}\bigr)
    \right\|_{\alpha_N,\lambda_N}
    \\
    &\qquad
    \le
    \mathcal R_{M_N}^{\alpha_N,\beta_N,\lambda_N}
    +
    \frac{
    2B_f
    }{
    d_N(2\pi)^{d_N}
    }
    N^{-d_N}
    \mathcal K_{M_N,d_N}^{\alpha_N,\beta_N,\lambda_N}.
\end{aligned}
}
\]
This identity shows explicitly that the truncation term is controlled by
\[
    N^{-d_N}
    \mathcal K_{M_N,d_N}^{\alpha_N,\beta_N,\lambda_N},
\]
where
\[
    d_N<1+\alpha_N
    \quad\text{and}\quad
    d_N<
    1+\frac{\beta_N+\lambda_N}{2}.
\]
Hence increasing $\alpha_N$ improves the vanishing of the testing kernel at
$x=1$, while increasing $(\beta_N+\lambda_N)/2$ improves the vanishing at
$x=0$. The singularity-matching condition is
\[
    \frac{\beta_N-\lambda_N}{2}\approx r.
\]
A natural parameterization is therefore
\[
    \lambda_N=\ell N+\lambda_0,
    \qquad
    \beta_N=\lambda_N+2r,
    \qquad
    \alpha_N=aN,
    \qquad
    M_N=\sigma N,
\]
with positive constants $a,\ell,\sigma$ and with
\[
    d_N=\lfloor \theta N\rfloor,
    \qquad
    0<\theta<\min\{a,\ell\}.
\]
The remaining nontrivial task is to prove that, under this scaling,
\[
    \frac{
    2B_f
    }{
    d_N(2\pi)^{d_N}
    }
    N^{-d_N}
    \mathcal K_{M_N,d_N}^{\alpha_N,\beta_N,\lambda_N}
\]
does not grow faster than the damping factor
$N^{-d_N}(2\pi)^{-d_N}$. This is the GLOF analogue of the
Fourier--Gegenbauer truncation-error estimate.

%------------------------------------------------------------
\begin{remark}[About standard unweighted $L^2(0,1)$]
If one takes $\alpha=\lambda=0$, then
\[
    \|\cdot\|_{\alpha,\lambda}=\|\cdot\|_{L^2(0,1)}.
\]
In this case the projection is orthogonal in the standard $L^2$ norm.
However, the endpoint-vanishing condition
\[
    d<1+\min\left\{\alpha,\frac{\beta+\lambda}{2}\right\}
\]
does not allow any positive integer $d$ solely through the above
boundary-vanishing argument, because $\alpha=0$ gives no vanishing at
$x=1$. Thus the present high-order integration-by-parts estimate is
naturally a weighted $L^2_{\chi_{\alpha,\lambda}}$ estimate with
$\alpha>0$. In the standard $L^2$ setting, one still has the stability
bound
\[
    \left\|
    \pi_M^{0,\beta,0}
    \bigl(f-f_N^{\rm F}\bigr)
    \right\|_{L^2(0,1)}
    \le
    \|f-f_N^{\rm F}\|_{L^2(0,1)},
\]
but this only gives the usual Fourier-tail rate and does not by itself
produce an exponential truncation estimate.
\end{remark}
\end{document}